\documentclass[12pt]{article}
% Préambule commun à tous les documents extraits de « La Revue CAPES »
\usepackage[T1]{fontenc}
\usepackage[utf8]{inputenc}
\usepackage{lmodern}
\usepackage{amsmath,amssymb,amsthm,amsfonts,mathrsfs}
\usepackage{setspace}
\usepackage{listings}
\usepackage{pgfplots}
\pgfplotsset{compat=1.18}
\usepackage{longtable}
\usepackage{adjustbox}
\usepackage{lettrine}
\usepackage{tkz-tab}
\usepackage{changepage}
\usepackage{pdflscape}
\usepackage{graphicx}
\usepackage{tikz}
\usepackage{xcolor}
\usepackage[a4paper,top=2cm,bottom=2.5cm,left=2.5cm,right=2cm]{geometry}
\usepackage{fancyhdr}
\usepackage{draftwatermark}
\SetWatermarkText{La RevueCapes}
\SetWatermarkLightness{0.9}
\SetWatermarkScale{2}
\usepackage{hyperref}
\hypersetup{colorlinks=true,linkcolor=blue!50!black,urlcolor=blue!50!black}

\definecolor{revue}{RGB}{20,60,120}

\pagestyle{fancy}
\fancyhf{}
\renewcommand{\headrulewidth}{0.4pt}
\setlength{\headheight}{15pt}

% \entete{Matière}{Titre}{Type de document}
\newcommand{\entete}[3]{%
  \fancyhead[L]{\small\textsc{La Revue CAPES} -- #1}%
  \fancyhead[R]{\small #2}%
  \fancyfoot[C]{\thepage}%
  \thispagestyle{plain}%
  \begin{center}
    {\color{revue}\rule{\textwidth}{1.2pt}}\\[4pt]
    {\small\textsc{Concours directs CAP-CEG / CAPES -- Burkina Faso}}\\[6pt]
    {\LARGE\bfseries\color{revue} #2}\\[6pt]
    {\large\itshape #1 -- #3}\\[2pt]
    {\color{revue}\rule{\textwidth}{1.2pt}}
  \end{center}
  \vspace{0.5cm}%
}

% Encadré pour les remarques ajoutées dans les corrigés
\newcommand{\remarque}[1]{\par\medskip\noindent\fcolorbox{revue}{revue!6}{\parbox{\dimexpr\linewidth-2\fboxsep-2\fboxrule}{\textbf{\color{revue}Remarque.} #1}}\par\medskip}


\begin{document}
\entete{Mathématiques}{Sujet type n°7}{Énoncé}
\begin{spacing}{1.5}

 
 \medskip\textbf{Exercice 1}

Un point M, de coordonnées $x$ et $y$, décrit une courbe $C$. Les coordonnées $x$ et $y$ de M 
dépendent d'un paramètre réel $t$, selon les relations suivantes : 


\[x(t) =\frac{t^2}{t^2+1}\]

\[y(t) =\frac{t^3}{t^2+1}
\]

1) Quel est l'ensemble des valeurs possibles pour $x$ ? pour $y$ ? 

2) Montrer que le point $O(0,0)$ appartient à la courbe $C$.  

On rappelle que la pente $p$ de la tangente en un point de coordonnées $(x_0,y_0)$ est telle que 
$(y - y_0)/(x - x_0) = p$. Quelle est la pente de la tangente à $C$ au point $O$ ? 

3) Donner l'équation de la courbe $C$ en coordonnées cartésiennes sous la forme $y = g(x)$. 

Calculer $y^\prime$ et donner les variations de $g$ et la forme générale de la courbe $C$.

\medskip\textbf{Exercice 2}

 Soit $f$ la fonction définie sur $\mathbb{R}_{+}^{*}$ par
 \[
 \begin{cases}
 f(0)=0\\
 f(x)= \frac{x^2}{sh x}
 \end{cases}
 \]
 
 1. Déterminer le développement limité d'ordre 5 de $e^x$ et de $sh x$ en 0.
 
 2. Déduire le développement limité d'ordre 3 de $f$ en 0.
 
 3. En déduire que $f$ est dérivable sur $\mathbb{R}$ et préciser la position relative au voisinage de 0 du graphe de $f$ par rapport à sa tangente en 0.
 
 4. En déduire $f^{(3)}(0)$.
 
 \medskip\textbf{Exercice 3}

 Un cadenas possède un code à 3 chiffres, chacun des chiffres pouvant être un chiffre de 1 à 9.
 
 1. Dans cette question, les trois chiffres du code sont quelconques.
 
 (a) Combien y-a-t-il de codes possibles?
 
 (b) Déterminer le nombre de codes se terminant par un chiffre impair. En déduire le nombre de codes qui se terminent par un chiffre pair,
 
 (c) Combien y-a-t-il de codes contenant au plus un chiffre 5?
 
 (d) Combien y-a-t-il de codes ne contenant pas le chiffre 5.
 
 2. Dans cette question on souhaite que le code comporte obligatoirement trois chiffres distincts.
 
 (a) Combien y-a-t-il de codes possibles ?
 
 (b) Déterminer le nombre de codes se terminant par un chiffre pair. En déduire le nombre de codes se terminant par un chiffre impair.
 
 (c) Combien y-a-t-il de codes comprenant les chiffres 5 et 7?
 
 \medskip\textbf{Exercice 4}

 Soit $f$ la fonction ainsi définie sur $\mathbb{R}^2$ par:
 
 \[f(x,y)=2x^3+xy^2+5x^2+y^2\]
 
 1. La fonction $f$ possède-t-elle un maximum,ou un minimum global dans $\mathbb{R}^2$.
 
 2. Déterminer les points ou $f$ présente des extremum locaux dans $\mathbb{R}^2$.
 
 
 \quad

\end{spacing}
\end{document}
